% ECONOMICS_PROBLEM: {"id": "D43-359", "title": "Competition between sellers of positional goods", "jel_code": "D43", "source_id": "OP-0185"}
% !TeX program = xelatex
\documentclass[11pt,a4paper]{article}
\usepackage[margin=23mm]{geometry}
\usepackage{fontspec}
\setmainfont{Latin Modern Roman}
\usepackage{amsmath,amssymb,mathtools}
\usepackage{xcolor,enumitem,booktabs,longtable,array,xurl,hyperref}
\definecolor{muted}{HTML}{53636C}
\hypersetup{colorlinks=true,urlcolor=blue,linkcolor=blue}
\setlength{\parindent}{0pt}
\setlength{\parskip}{5pt}
\newcommand{\R}{\mathbb R}
\newcommand{\E}{\mathbb E}
\newcommand{\N}{\mathbb N}
\newcommand{\ind}{\mathbf 1}
\newcommand{\Var}{\operatorname{Var}}
\newcommand{\Cov}{\operatorname{Cov}}
\newcommand{\supp}{\operatorname{supp}}
\DeclareMathOperator*{\argmin}{arg\,min}
\DeclareMathOperator*{\argmax}{arg\,max}
\newcommand{\entrymeta}[3]{{\sffamily\small\color{muted}\textbf{\texttt{#1}}\quad|\quad #2\par #3\par}}
\newcommand{\Problemref}[1]{\texttt{#1}}
\newcommand{\JELref}[2]{\texttt{#2} (JEL \texttt{#1})}
\begin{document}
\section*{Competition between sellers of positional goods}
\textbf{Problem ID:} \texttt{D43-359}\quad \textbf{JEL:} \texttt{D43}\par
% BEGIN ECONOMICS STATEMENT
\label{problem:OP-0185}
% GAME_THEORY_PROBLEM: {"local_id": "GT-055", "original_number": 55, "kind": "R", "entry_kind": "statement_only", "published": false, "review_required": true, "category": "game_theory"}
\label{game:GT-055}
{\sffamily\small\color{muted}
\textbf{GT-055} | Mechanism competition and externalities | \textbf{R}: explicitly specified finite-menu duopoly.
\par}
\subsection*{Definitions and mathematical statement}
A unit mass of consumers has types $\theta\in[0,1]$ with a positive continuous density $f$. Two sellers $s\in\{1,2\}$ each offer fixed quality levels $0<q_{s1}<\cdots<q_{sK}$ and choose prices $p_{sk}\in[0,B]$. Fix continuous nonnegative intrinsic utilities $a_{sk}(\theta)$ and costs $c_{sk}$. A population response is a measurable probability vector $z(\theta)$ over the outside option and all seller-level pairs. Let $m_{sk}=\int_0^1z_{sk}(\theta)f(\theta)d\theta$. Status within seller $s$ is
\[
 S_{sk}(m)=\sum_{\ell<k}m_{s\ell}+\frac12m_{sk}.
\]
Consumers assigned to $(s,k)$ obtain $a_{sk}(\theta)+\theta S_{sk}(m)-p_{sk}$; outside utility is zero. A response equilibrium at prices $p$ satisfies: $z_{sk}(\theta)>0$ only if $(s,k)$ maximizes these utilities and zero, for almost every $\theta$. Seller profit is $\Pi_s(p,z)=\sum_k(p_{sk}-c_{sk})m_{sk}$.
Characterize the inputs for which there exist a response-equilibrium selection $z(p)$ for every price vector and prices $p^*$ satisfying
\[
 \Pi_s(p^*,z(p^*))\geq
 \Pi_s(p'_s,p^*_{-s},z(p'_s,p^*_{-s}))
 \quad\forall s,\ p'_s\in[0,B]^K.
\]
The target includes the resulting allocation/welfare correspondence and its comparison with the monopolist benchmark under the same intrinsic utilities and costs.
\subsection*{Short English statement}
Characterize credible seller price competition when status depends on the customers of each seller.
\subsection*{Variants and scope boundaries}
Menus and the two-stage selection rule are fixed explicitly. Allowing sellers to choose quality distributions or arbitrary direct mechanisms broadens the agenda. Across-seller status comparisons and normalization by each seller's customer mass produce different externalities.
\subsection*{Source relationship and status}
The source explicitly leaves strategic mechanism competition unexplored but specifies no canonical game or conjecture. This finite-menu model is editorial and narrows the agenda to a precise equilibrium-existence/classification target. Its status is therefore R, not a published claim that all such games lack known equilibria.
\subsection*{References}
\begingroup\footnotesize\raggedright
\textbf{[S1]} Peiran Xiao (2025 version, first posted 2024). \emph{Allocating Positional Goods: A Mechanism Design Approach}. arXiv:2411.06285v3. \textit{Verified locator:} Section 2.1, status definition; Section 5, concluding two paragraphs on seller competition and separate queues. \url{https://arxiv.org/html/2411.06285v3}
\par\endgroup

\subsection*{Common conventions: Source mathematical conventions}

\textbf{Finite games.} For a finite set $A$, $\Delta(A)$ is its probability
simplex. Players choose independent mixed strategies in Nash equilibrium;
correlation is allowed only when explicitly part of the solution concept.
In a game with payoff functions $u_i$ and strategy profile $x$, an additive
$\varepsilon$ Nash equilibrium satisfies
\[
 u_i(x)\geq u_i(x_i',x_{-i})-\varepsilon
 \quad\text{for every player $i$ and every admissible deviation $x_i'$}.
\]
For numerical approximation, payoffs are normalized as locally specified.
An $\varepsilon$ payoff residual is distinct from distance at most
$\varepsilon$ to the set of exact equilibria. Point convergence, convergence
of empirical play, and convergence in distance to an equilibrium set are
also distinct.

\textbf{Computational representations.} Unless stated otherwise, a complexity
question takes rational data with binary encoding; polynomial time is measured
in total bit length. A fixed-accuracy polynomial algorithm is not necessarily
polynomial in $1/\varepsilon$, and neither is necessarily polynomial in
$\log(1/\varepsilon)$. Succinct games and value, demand, or sampling oracles
must declare their interfaces. Exact real solutions require an explicit
representation; an unspecified list of arbitrary real numbers is not a
finite computational output. A claimed algorithmic barrier is meaningful
only for its specified model and reductions.

\textbf{Dynamic games.} Histories, information, observation, and allowable
randomization are part of the model. A stationary strategy depends only on
the current state; a positional strategy in a deterministic graph game
chooses one action at each controlled vertex. Nash equilibrium at an initial
state and equilibrium after every history are different requirements.
An expectation of a pathwise $\liminf$ payoff, a $\liminf$ of expected averages,
a limiting expected average, and a uniform finite-horizon equilibrium are
different criteria. None is silently substituted for another.

\textbf{Allocation and mechanisms.} A complete allocation assigns every item
exactly once unless otherwise stated. Zero-valued goods, negative values,
capacity constraints, feasibility, lotteries, and copies of goods affect
fairness definitions and are specified locally. Dominant-strategy incentive
compatibility, Bayesian incentive compatibility, universal truthfulness,
and truthfulness in expectation impose different inequalities. Whenever
an objective ratio has zero denominator, its convention is stated or those
instances are excluded.

\textbf{Infinite-dimensional models.} Expectations, integrals, stochastic
equations, and optimization objectives require their local regularity,
integrability, filtrations, and admissible domains. A backward--forward
mean-field system is a boundary-value problem; it is not an initial-value
semiflow without an additional construction. Long-time, large-population,
and vanishing-noise limits require an order of limits and a convergence
topology. If a minimum need not be attained, the objective is its infimum
or supremum and the approximation requirement is stated separately.
% END ECONOMICS STATEMENT
\end{document}
